\documentclass[10pt,a4paper]{amsart}
\usepackage[margin=19mm]{geometry}
\usepackage{amsmath,amssymb,mathtools}
\usepackage[scr=rsfs,cal=euler]{mathalfa}
\usepackage{xfrac}
\usepackage[unicode,hidelinks,bookmarks=false]{hyperref}
\newcommand{\ie}{i.e.\ }
\newcommand{\cf}{cf.\ }
\newcommand{\resp}{respectively\ }
\newcommand{\Subsection}{\subsection}
\providecommand{\nobreakdash}{\nobreak\hbox{-}}
\setlength{\emergencystretch}{2em}
\multlinegap0pt
\usepackage{bm}
\usepackage{bbm}
\usepackage{dsfont}
\usepackage{xspace}
\usepackage{cancel}
\usepackage{mathtools}
\usepackage{cleveref}
\usepackage{amsthm}
\makeatletter
\@ifundefined{convention}{\newtheorem{convention}{Convention}}{}
\@ifundefined{theorem}{\newtheorem{theorem}{Theorem}}{}
\@ifundefined{definition}{\newtheorem{definition}[theorem]{Definition}}{}
\@ifundefined{remark}{\newtheorem{remark}[theorem]{Remark}}{}
\makeatother
\DeclareMathOperator{\supp}{supp}
\newcommand\ddfrac[2]{\frac{\displaystyle #1}{\displaystyle #2}}
\newcommand{\RQ}{\textrm{RQ}}
\title[Coloring outside the lines: Spectral bounds for generalized ]{Coloring outside the lines: Spectral bounds for generalized hypergraph colorings}
\author{Lies Beers, Raffaella Mulas}
\date{Source: arXiv:2506.17659v2 (CC BY 4.0)}
\begin{document}
\maketitle

\noindent\emph{Source and licence.} Coloring outside the lines: Spectral bounds for generalized hypergraph colorings. Version arXiv:2506.17659v2 (CC BY 4.0), distributed under the Creative Commons Attribution 4.0 International licence, as recorded in the arXiv licence metadata for this exact version and shown on the abstract page https://arxiv.org/abs/2506.17659v2. Full rights notice: licenses/arxiv-2506.17659-CC-BY-4.0.txt.\par\smallskip

\noindent\textit{Editorial scope.} Counted unit: the Convention whose source label is \texttt{conv:1} in the article (source0.tex lines 427--433, source label \texttt{conv:1}), delivered together with the article's own complete original proof of it (source0.tex lines 435--541, one \texttt{proof} environment of 107 lines). No mathematics is written, repaired or re-derived: statement and proof are the article's own text line for line. Required context retained uncounted: rmk:A=0 (lines 276--285); def:coloring1 (lines 372--374); thm:cor5.4AMZ (lines 418--423). Every definition, display and label of those lines is retained unchanged. Omitted from this package and not redistributed: 1--275, 286--371, 375--417, 424--426, 434--434, 542--1334. Nothing inside a retained excerpt is omitted or altered: every retained body is delivered exactly as the article's lines are written, so no omission receipt of any kind accompanies this package. Citations: the retained excerpts carry no citation command at all, so no bibliography entry of the article is transcribed and no reference is redistributed. Reference adaptation: none was needed and none was made; every reference inside the delivered unit resolves inside it, so no unresolved reference or citation is produced. Printed slips kept literal: no printed word, symbol, hypothesis or number of the counted statement or of its proof was altered. Layout and class: the article's own class is \texttt{article} with packages including inputenc, amsmath, amssymb, amsthm, amsfonts, latexsym, graphicx, subfigure, mathtools, enumitem, url, xcolor, svg, stmaryrd; the wrapper is the lane's pinned \texttt{amsart} editorial wrapper at 10pt on A4, which restates the theorem-like environments the retained text needs and adds the article's own macro definitions that the retained text needs (\texttt{\textbackslash RQ}, \texttt{\textbackslash ddfrac}, \texttt{\textbackslash supp}), with extra wrapper packages (bm, bbm, dsfont, xspace, cancel, mathtools, cleveref). The delivered numbering is the unit's own. Assets: no retained span contains a figure environment, a graphics-inclusion command, a PDF-page inclusion, a table or a TikZ picture; no image file is copied into this package, the package declares no asset dependency and no image file is redistributed with it. Licence: version arXiv:2506.17659v2 of this article is distributed under the Creative Commons Attribution 4.0 International licence, as recorded in the arXiv licence metadata for this exact version and shown on the abstract page https://arxiv.org/abs/2506.17659v2. A full rights notice is delivered at licenses/arxiv-2506.17659-CC-BY-4.0.txt. Characters: every retained excerpt is pure ASCII, so the delivered engine, XeLaTeX with its default Latin Modern text font, reports no missing character.
\begin{remark}\label{rmk:A=0}
We have 
\begin{align*}
    A=\textbf{0} &\iff L=I \\
    &\iff \lambda_1=\ldots=\lambda_N=1\\
    &\iff \lambda_1=1\\
    &\iff \lambda_N=1.
\end{align*}
The last two characterizations follow from the fact that $\sum_{i=1}^N\lambda_i=N$. Hence, if $$1=\lambda_1\leq \lambda_2\leq \ldots\leq\lambda_N,$$ then $1=\lambda_1= \ldots=\lambda_N$. Similarly, if $$\lambda_1\leq \lambda_2\leq \ldots\leq\lambda_N=1,$$ then $\lambda_1= \ldots=\lambda_N=1$. In particular, this implies that, if $A\neq \textbf{0}$, then $\lambda_1<1$ and $\lambda_N>1$.
\end{remark}

\begin{definition}\label{def:coloring1}
    Let $\Gamma$ be an oriented hypergraph. A \emph{proper strong (vertex) $k$-coloring} is a function $c\colon V(\Gamma)\to \{1,\ldots,k\}$ such that, for every edge $e$ and any two vertices $v_1,v_2\in e$, we have $c(v_1)\neq c(v_2)$. The corresponding coloring number is denoted by $\chi\coloneqq \chi(\Gamma)$.
\end{definition}

\begin{theorem}\label{thm:cor5.4AMZ}
    Let $\Gamma$ be an oriented hypergraph and let $\chi(\Gamma)$ be defined as in Definition \ref{def:coloring1}. Then, we have that
    \[
    \chi(\Gamma) \geq \frac{\lambda_N - \lambda_1}{\min\bigl\{ \lambda_N-1, 1-\lambda_1 \bigr\}}.
    \]
\end{theorem}

\begin{convention}\label{conv:1}
We recall that, by Remark \ref{rmk:A=0}, we have $\lambda_1=1$ if and only if $\lambda_N=1$. We use the convention that, in this case, 

\[
\frac{\lambda_N-\lambda_1}{\lambda_N-1} = \frac{\lambda_N-\lambda_1}{1-\lambda_1} = 1.
\]
\end{convention}

\begin{proof}[Alternative proof of Theorem \ref{thm:cor5.4AMZ}]
    We first consider the case where $\lambda_1 = \lambda_N = 1$. By Convention \ref{conv:1}, in this case, the inequality in the statement becomes
    \[
    \chi(\Gamma)\geq 1,
    \]
    which is trivially true.

    Now assume that $\lambda_1<1$ and hence, by Remark \ref{rmk:A=0}, $\lambda_N > 1$.
    We first show that $$\chi(\Gamma)\geq \frac{\lambda_N-\lambda_1}{1-\lambda_1}.$$ Fix a proper $\chi$-coloring, and let $V_1,\ldots,V_\chi$ be the corresponding coloring classes. Furthermore, let $g$ be an eigenfunction for $\lambda_N$, so that $\RQ(g)=\lambda_N$.
    Let also $l\leq \chi$ be such that
    \[
    \supp(g)\subseteq \bigcup_{i=1}^lV_i,
    \]
    where we rearrange the coloring classes if necessary, to ensure that $l$ is as small as possible. Note that $l\geq2$, because, if $\supp(g)\subseteq V_1$, then $\RQ(g)= 1$.

    Now define, for $1\leq i<j\leq  l$,
    \begin{equation}\label{eq:S_ij}
    S_{ij}^g\coloneqq  \sum_{v_i\in V_i, v_j\in V_j}A_{v_i,v_j}g(v_i)g(v_j).
    \end{equation}
    We have that
    \begin{align*}
        \lambda_N &= \RQ(g) \\&= 1 - 2\ddfrac{\sum_{v,w\in V}A_{v,w}\cdot g(v)\cdot g(w)}{\sum_{v\in V}\deg v\cdot g(v)^2}\\
        &= 1 - 2\ddfrac{\sum_{1\leq i<j\leq  l}\sum_{v_i\in V_i,v_j\in V_j}A_{v_i,v_j}\cdot g(v_i)\cdot g(v_j)}{\frac1{ l-1}\sum_{1\leq i< j\leq  l}\sum_{v\in V_i\cup V_j}\deg v \cdot g(v)^2}\\
        &\leq 1 + 2\ddfrac{\sum_{1\leq i<j\leq  l}\max\biggl\{-S^g_{ij},0\biggr\}}{\frac1{ l-1}\sum_{1\leq i< j\leq  l}\sum_{v\in V_i\cup V_j}\deg v \cdot g(v)^2}.
    \end{align*}
    Note that, since  $\lambda_N>1$, we must have that $S^g_{ij}< 0$ for some $i,j$. Now we let, for all $i,j$,
     \begin{equation}\label{eq:g_ij}
     g_{ij}(v)\coloneqq \begin{cases}
        g(v), &\text{ if } v\in V_i,\\
        -g(v), &\text{ if } v\in V_j,\\
        0, &\text{ otherwise.}
    \end{cases}
    \end{equation}
    and we assume without loss of generality that
    \[
    \RQ(g_{12}) = \min_{1\leq i<j\leq  l} \RQ( g_{ij}).
    \]
    Note that $S_{12}^g<0$.
    Then,
    \begin{align*}
        \lambda_N &\leq 1 + 2\ddfrac{\sum_{1\leq i<j\leq  l}\max\biggl\{-S^g_{ij},0\biggr\}}{\frac1{ l-1}\sum_{1\leq i< j\leq  l}\sum_{v\in V_i\cup V_j}\deg v \cdot g(v)^2}\\
        &\leq 1 + 2( l-1)\ddfrac{\bigl|S^g_{12}\bigr|}{\sum_{v\in V_1\cup V_2}\deg v \cdot g(v)^2}\\
        &\overset{(i)}{=} 1+\bigl( l-1\bigr)\bigl(1-\RQ(g_{12})\bigr).
    \end{align*}
    To prove $(i)$, recall that $S_{12}^g<0$, and note that
    \begin{align*}
        \RQ\bigl(g_{12}\bigr) &= 1 -2\frac{\sum_{\substack{v\in V_1\\ w\in V_2}}A_{v,w}\cdot g(v)\cdot -g(w)}{\sum_{v\in V_1\cup V_2}\deg v\cdot g(v)^2}\\
        &= 1-2\frac{\bigl|S_{12}^g\bigr|}{\sum_{v\in V_1\cup V_2}\deg v\cdot g(v)^2}.
    \end{align*}
    From the equality $\lambda_N\leq 1+(l-1)(1-\RQ(g_{12}))$, it follows that
    \begin{equation*}
        \lambda_N \leq 1+\bigl( l-1\bigr)\bigl(1-\lambda_1\bigr),
    \end{equation*}
     which is equivalent to
    \begin{equation}\label{eq:inequality1.1}
        \frac{\lambda_N-\lambda_1}{1-\lambda_1} \leq l \leq \chi.
    \end{equation}
    The proof of the inequality
    \[
    \chi(\Gamma) \geq \frac{\lambda_N-\lambda_1}{\lambda_N-1}
    \]
    is analogous. In particular, let $h$ be an eigenfunction such that $\RQ(h) = \lambda_1$, and let $k\leq \chi$ be such that
    \[
    \supp(h)\subseteq \bigcup_{i=1}^kV_i,
    \]
    where we rearrange the coloring classes if necessary, to ensure that $k$ is as small as possible. Note that $k\geq2$, because, if $\supp(h)\subseteq V_1$, then $\RQ(h)= 1$.

    
    We define, for all $1\leq i<j\leq k$, the functions
    \begin{equation}\label{eq:h_ij}
     h_{ij}(v)\coloneqq \begin{cases}
        h(v), &\text{ if } v\in V_i,\\
        - h(v), &\text{ if } v\in V_j,\\
        0, &\text{ otherwise.}
    \end{cases}
    \end{equation}
    We assume without loss of generality that
    \[
    \RQ( h_{12}) = \max_{1\leq i<j\leq k}\RQ(h_{ij}),
    \]
    and we let $S^h_{ij}$ be defined as in \cref{eq:S_ij}. We see that
    \begin{align*}
        \lambda_1 &= \RQ(h) \\&= 1 - 2\ddfrac{\sum_{v,w\in V}A_{v,w}\cdot h(v)\cdot h(w)}{\sum_{v\in V}\deg v\cdot h(v)^2}\\
        &= 1 - 2\ddfrac{\sum_{1\leq i<j\leq  k}\sum_{v_i\in V_i,v_j\in V_j}A_{v_i,v_j}\cdot h(v_i)\cdot h(v_j)}{\frac1{ k-1}\sum_{1\leq i< j\leq  k}\sum_{v\in V_i\cup V_j}\deg v \cdot h(v)^2}\\
        &\geq 1 - 2\ddfrac{\sum_{1\leq i<j\leq  k}\max\bigl\{S^h_{ij},0\bigr\}}{\frac1{ k-1}\sum_{1\leq i< j\leq  k}\sum_{v\in V_i\cup V_j}\deg v \cdot h(v)^2}\\
        &\geq 1 - 2( k-1)\ddfrac{S_{12}^h}{\sum_{v\in V_1\cup V_2}\deg v \cdot h(v)^2}\\
        &\overset{(ii)}{=} 1-\bigl( k-1\bigr)\bigl(\RQ(h_{12})-1\bigr).
    \end{align*}     
        
    To prove $(ii)$, we use the fact that $S_{12}^h>0$ (since $\RQ(h)<1$), and that
    \begin{align*}
        \RQ\bigl(h_{12}\bigr) &= 1 -2\frac{\sum_{\substack{v\in V_1\\ w\in V_2}}A_{v,w}\cdot h(v)\cdot -h(w)}{\sum_{v\in V_1\cup V_2}\deg v\cdot h(v)^2}\\
        &= 1+2\frac{S_{12}^h}{\sum_{v\in V_1\cup V_2}\deg v\cdot h(v)^2}.
    \end{align*}
    By assumption, it now follows that
    \begin{equation*}
    \begin{aligned}
        \lambda_1&\geq 1-( k-1)(\lambda_N-1) \\
        &=  k(1-\lambda_N)+\lambda_N,
    \end{aligned}
    \end{equation*}
    i.e.,
    \begin{equation}\label{eq:inequality2.2}
        \chi \geq k \geq \frac{\lambda_N-\lambda_1}{\lambda_N-1}.
    \end{equation}
    This concludes the proof.
\end{proof}

\end{document}
